\section{Introduction}
\label{sec:intro}

Whether uninsured depositors discipline bank risk remains one of the
central open questions in banking.  Theory predicts they should
\citep{calomiris1991role, diamond2001liquidity}, and a large empirical
literature documents that riskier banks pay higher deposit rates and
experience larger uninsured outflows \citep{park1995market,
martinez2001depositors, maechler2006dynamic, iyer2016tale, egan2017deposit,
chen2024liquidity}.  Whether these patterns reflect a causal disciplinary
response or equilibrium correlation remains unsettled.  The question is
pressing: against a backdrop of expanding government protection for
uninsured depositors---the Main Street Depositor Protection Act (S.~2999,
reintroduced March 2026) would formally extend coverage to \$10~million
per account---it matters whether this discipline operates before the
incentives that sustain it are further weakened.  Early state deposit
insurance programs eliminated market discipline and enabled excessive
risk-taking \citep{calomiris2019stealing}.  We ask whether uninsured
depositors respond to mandatory credit-risk disclosures, using
U.S.\ community banks, where institutions below \$10~billion in assets
carry no credible government backstop and identification conditions are
favorable \citep{o1990deposit, flannery1996evidence}.

Clean identification requires resolving three problems.  Simultaneity: a
bank's risk profile, deposit composition, and depositor behavior
co-evolve, so regressions of flows on risk measures capture equilibrium
sorting rather than isolated disciplinary responses.  Accounting
opacity: under the incurred-loss model, banks could delay recognition
and systematically understate risk in disclosures
\citep{bushman2012accounting, bushman2015delayed}.  Aggregate stress:
periods with the greatest variation in bank distress are precisely those
in which government interventions and depositors' own liquidity needs
contaminate the response \citep{bennett2015market}.  We address the
first two by exploiting mandatory CECL adoption for community banks in
2023 as an information shock whose cross-sectional intensity was fixed
before the event.  CECL requires banks to recognize lifetime expected
credit losses on existing loan portfolios at adoption, producing a
one-time adjustment to retained earnings whose magnitude reflects loan
vintages assembled years earlier.  Banks knew the standard was coming,
but the size of the Day-One charge---the treatment variation we
use---was pinned down by the composition and seasoning of the existing
loan book rather than by contemporaneous deposit-market conditions.
The third problem we confront directly: community banks adopted CECL in
the same quarter that the failures of Silicon Valley Bank and Signature
Bank put bank risk on every depositor's mind.  We show that our
estimates are robust to allowing deposit dynamics to vary flexibly with
the full 2022Q4 balance-sheet characteristics that governed exposure to
the 2023 stress---uninsured deposit shares, size, capital, asset
quality, profitability, CRE concentration, and unrealized securities
losses---each interacted with calendar-quarter fixed effects, and that
response patterns line up with each cohort's own adoption quarter rather
than with March 2023.  \citet{kim2026current} and \citet{gee2022decision}
confirm the Day-One adjustment is informative;
\citet{nier2006market} and \citet{chen2022bank} establish that
disclosure improvements enable depositors to act on credit-risk signals.

Our central result, estimated on more than 4,000 community banks over
quarterly data from 2016 to 2025, is a within-bank shift in the
composition of time deposits away from uninsured and toward insured
balances.  A triple-difference specification with bank-by-quarter fixed
effects implies that each percentage point of CECL exposure reduces
uninsured time deposits by 0.18 percentage points of assets
\emph{relative to} insured time deposits within the same bank and
quarter.  This composition result survives characteristics-by-quarter
controls, alternative treatment scalings and cutoffs, and the exclusion
of the largest-adjustment bank.  The level DiD is consistent: banks in
the right tail of CECL exposure lost 0.394 percentage points of
uninsured time deposits (roughly 7 percent of the sample mean) and
gained 0.712 percentage points of insured time deposits.  Event studies
confirm parallel pre-adoption trajectories and a divergence that emerged
at adoption and persisted.  The level uninsured decline attenuates when
deposit-composition characteristics receive their own time paths, as it
mechanically must, since those controls absorb part of the uninsured
channel itself---which is why we anchor the paper on the within-bank
composition contrast.

High-CECL banks replace the lost uninsured balances with \emph{insured}
wholesale funding: insured brokered deposits rise by roughly 0.5
percentage points of assets and reciprocal deposits---which split large
accounts across network banks to obtain full coverage---also increase,
while uninsured brokered funding does not move.  Branch-level SOD data
show that deposit growth after 2023 is faster at low-CECL banks
operating in counties more exposed to high-CECL competitors, direct
evidence that part of the withdrawn funding re-intermediates locally.
Both margins are what a disciplinary mechanism predicts.  The
reallocation carries real costs.  High-CECL banks paid approximately
8--9 basis points more annually in interest expense, ROA fell by
roughly 15 basis points, and NIM compressed by approximately 9 basis
points.  Quarterly loan growth slows by 0.65 percentage points,
concentrated in CRE, linking the funding shock to credit supply in the
spirit of \citet{granja2025current}.  Effects are amplified by depositor
sophistication: high-CECL banks in high-sophistication markets paid an
additional 17--18 basis points annually in interest expense and
suffered an additional 30 basis point drag on ROA; the pattern is
unchanged under a pre-period (2019) branch-weight index.

This paper contributes to three literatures.  First, on depositor and
market discipline, prior evidence documents higher rates and larger
outflows at riskier banks \citep{park1995market, egan2017deposit,
flannery1996evidence, nier2006market}, but causal interpretation has
been elusive: government guarantees attenuate estimates in the crisis
episodes that generate the most variation \citep{bennett2015market,
berger2015depositors, jacewitz2018deposit}, and cross-sectional
comparisons cannot separate information responses from equilibrium
sorting.  Recent work on the 2023 episode shows depositors ran on banks
with high uninsured shares and unrealized losses
\citep{caglio2024depositors}; we exploit orthogonal variation
\emph{within} that high-attention period, uncorrelated with the
balance-sheet characteristics that governed run exposure.  Second, on
CECL and mandatory credit-risk disclosure, \citet{bushman2012accounting}
and \citet{bushman2015delayed} document that delayed loss recognition
under the incurred-loss model attenuated the informativeness of bank
statements and, by extension, the market's ability to price bank risk.
Prior CECL work focuses on the supply side---lending
\citep{granja2025current}, provisioning \citep{kim2026current}---and on
information users other than depositors, particularly analysts
\citep{gee2022decision} and equity investors.  We complement this
literature by tracing the disclosure to depositor behavior and
downstream to funding costs, profitability, and credit supply, closing
a channel that the analyst-and-equity strand does not observe.  Third,
on deposit insurance design, broader coverage reduces the incentive to
monitor \citep{demirgucc2004market, flannery2019market}; the substitution
we document toward insured brokered and reciprocal deposits
\citep{kim2024insurance} shows that much of the response takes the form
of repurchasing insurance through market arrangements rather than
exiting the banking system.
